\section{Institutional Framework and Transaction Costs}
\label{sec:microstructure}

Running a quantitative arbitrage book in Mumbai is fundamentally different from running one in New York or London. Because retail and institutional short selling in the cash equity segment is held back by an illiquid borrow mechanism, quant managers seeking market neutrality almost always turn to single-stock futures on the National Stock Exchange.

\subsection{Statutory Indian Derivative Taxes}
In Western markets, institutional commissions are often fractions of a basis point, and high-volume firms even earn maker-taker rebates. In India, however, government taxes and exchange fees are fixed by statute and apply equally to all participants \citep{bose2007securities}:
\begin{enumerate}
    \item \textbf{Securities Transaction Tax (STT):} Levied strictly on the sell side of futures trades. From 2010 through September 2024, this was 0.0125\% ($1.25$ bps). On October 1, 2024, SEBI raised it by 60\% to 0.0200\% ($2.00$ bps) to curb speculative retail trading in derivatives.
    \item \textbf{Stamp Duty:} Imposed nationwide under the Indian Stamp Act at 0.0020\% ($0.20$ bps) on the buy leg.
    \item \textbf{Exchange Turnover Fees (NSE):} Blended charge of 0.0019\% ($0.19$ bps) on both buy and sell orders ($0.38$ bps round-trip).
    \item \textbf{SEBI Regulatory Fees:} A turnover fee of 0.0001\% ($0.01$ bps) on both sides ($0.02$ bps round-trip).
    \item \textbf{Goods and Services Tax (GST):} An 18.0\% tax levied on brokerage and exchange fees, adding about $0.12$ bps round-trip.
\end{enumerate}

\subsection{Slippage and Full Implementation Costs}
Beyond taxes, any real trading strategy faces execution slippage \citep{roll1984simple,lesmond1999new,amihud2002illiquidity,frazzini2018trading}. We model realistic slippage using Average Daily Traded Value (ADTV) tiers:
\begin{equation}
    \text{Slippage}_{\text{tier}} = \begin{cases}
        2.0 \text{ bps per leg (}4.0\text{ bps round-trip)} & \text{if ADTV} \ge \text{Rs.~100 Cr} \\
        5.0 \text{ bps per leg (}10.0\text{ bps round-trip)} & \text{if Rs.~25 Cr} \le \text{ADTV} < \text{Rs.~100 Cr} \\
        10.0 \text{ bps per leg (}20.0\text{ bps round-trip)} & \text{if ADTV} < \text{Rs.~25 Cr}
    \end{cases}
\end{equation}

Putting all components together, total round-trip cost $C_{i,t}$ for stock $i$ on date $t$ is:
\begin{equation}
    C_{i,t} = \text{STT}_t + \text{Stamp} + 2 \times (\text{NSE}_{\text{fee}} + \text{SEBI}_{\text{fee}}) + \text{GST} + 2 \times \text{Slippage}_i
\end{equation}
For an average large-cap stock, baseline round-trip friction was roughly $6.5$ basis points prior to October 2024, rising to $7.3$ basis points afterward. Crucially, a pair trade involves two distinct stocks. Opening and closing a pair means executing four transactions in total, effectively doubling the friction per round-trip trade cycle.
